\begin{lemma}[Proved: the $p$-th exterior Adams operation on every Jordan block]
\label{C:Jordan_Adams_p}
Let $k$ be a field of characteristic $p>2$, and let $V_j$ be the
indecomposable $k[C_p]$-module of dimension $j$, $1\leq j\leq p$.
In the split Green ring put $w=[V_p]-[V_{p-1}]$. Then $w^2=1$ and
\[
 \psi^p(V_j)=\begin{cases}j[V_1],&j\text{ odd},\\jw,&j\text{ even}.
 \end{cases}                                                   \tag{C.3.10}
\]
If $x=[V_2]$ and $D(a,b)=\psi^p(a)\psi^p(b)-\psi^p(ab)$, then
\[
 D(x,V_j)=0\quad(j<p),\qquad D(x,V_p)=2p(w-1).             \tag{C.3.11}
\]
\end{lemma}
\begin{proof}
The tensor calculations in the proofs of
Theorem~\ref{C:odd_characteristic_failure} and
Corollary~\ref{C:BO_positive_failure} give
$xV_j=V_{j-1}+V_{j+1}$ for $j<p$, $xV_p=2V_p$, and $w^2=1$.
These steps use actual Jordan kernels and the augmentation-ideal sequence;
they have been rechecked for their use here.
Define $U_0=0$, $U_1=1$, $U_{j+1}=xU_j-U_{j-1}$.
The first identity gives $V_j=U_j(x)$ for $j\leq p$.
Let $P_a$ be the integral polynomial determined by
$P_a(s+s^{-1})=s^a+s^{-a}$. Since $\lambda_t(V_2)=1+xt+t^2$,
Newton's recursion gives $\psi^a(x)=P_a(x)$ for every $a$.

For $a$ prime to $p$, the exterior Adams operation on a modular Green
ring is a ring endomorphism: Bryant--Johnson
\cite[Proposition 2.1]{C-BryantJohnson2011}. Applied with $1\leq a<p$,
this gives $\psi^a(V_j)=U_j(P_a(x))$.
In the torsion-free Green ring, Newton's recursions therefore determine
$\lambda^r(V_j)$ for $r<p$ to be the same integral polynomial in $x$
as the exterior character of the formal weights
$s^{j-1},s^{j-3},\ldots,s^{1-j}$. One may perform this comparison after
tensoring with $\mathbb Q$, where division by $1,\ldots,p-1$ is legitimate,
and descend because the Green ring has the free basis $V_1,\ldots,V_p$.
For $j<p$ these are every nonzero exterior power. For $j=p$ the one
remaining exterior power is the determinant $\lambda^p(V_p)=1$, also
equal to the product of those formal weights. Newton's formula at $p$
now gives
\[
 \psi^p(V_j)=U_j(P_p(x))=U_j(2w)=j w^{j-1}.
\]
Here $P_p(x)=2w$ was obtained directly from the recurrence at the boundary
$xV_p=2V_p$, and the last equality follows inductively from $w^2=1$.
This proves (C.3.10). Substitution in $xV_j=V_{j-1}+V_{j+1}$ shows
(C.3.11), for odd and even $j$ separately; at $j=p$ one instead uses
$xV_p=2V_p$ and $\psi^p(V_p)=p$.
\end{proof}

\begin{theorem}[Proved: the complete fieldwise special-$\lambda$ criterion]
\label{C:BO_field_lambda}
Let $k$ be any field of characteristic $p>2$, and let $O_m$ be the split
orthogonal group, including $O_0=1$. Then exterior powers make
$WG_{\mathrm{sp}}(B_{O_m})$ a special $\lambda$-ring if and only if
$m\leq2$. Thus the low-rank cases left by the bound $m\geq p+4$ are
all decided. The negative assertion holds over arbitrary perfect fields
and arbitrary algebraically closed fields, with no distinction in the
criterion.
\end{theorem}
\begin{proof}
For $m=0,1,2$ the representation category is semisimple. For $m=1$ this
is the eigenspace decomposition for $\mu_2$. For $m=2$ write
$O_2=\mathbb G_m\rtimes C_2$ with inversion action. Torus weights
split every representation; the zero weight splits into the two
$C_2$-characters, while each pair of nonzero opposite weights is the
sum of copies of the two-dimensional exchanged-weight representation.
This proves semisimplicity over $k$ itself. Every metabolic symmetric
representation then splits isometrically as a hyperbolic representation:
choose an equivariant complement to the Lagrangian and correct its
self-pairing by half its Gram map. Consequently the split group equals
the exact Grothendieck--Witt group. Zibrowius's main theorem, or his
Proposition 3.11 for $m=2$, applies to that group
\cite{C-Zibrowius2015}. Its field hypothesis and $2\ne0$ are exactly
our hypotheses, proving specialness.

We next treat $m\geq4$. Let $E_m$ be the standard symmetric
representation and give $B=\operatorname{Sym}^{p-1}(E_m)$ the symmetric
form obtained by restricting the tensor-power form to the image of the
symmetrizing idempotent. This idempotent exists since $(p-1)!$ is a
unit and is self-adjoint; its image is therefore nondegenerate.
Use the homomorphism $C_p\to O_m$ with representation
$H(V_2)\perp E_{m-4}^{\mathrm{triv}}$, where the second term is the
standard split quadratic space of dimension $m-4$ with trivial
$C_p$-action. It is defined over $k$: the hyperbolic form is split and
the indicated split orthogonal summand is fixed.
In the Green ring the restriction of $E_m$ is $2x+(m-4)$.
Write $b_j\geq0$ for the Jordan multiplicities of $B$ on this subgroup.
We have $b_p>0$: the direct-sum formula for symmetric powers exhibits
$\operatorname{Sym}^{p-1}(V_2)$ as a summand of the restriction, and
this module is $V_p$. Indeed on polynomials of degree at most $p-1$
the generator acts by $t\mapsto t+1$; its successive differences have
nonzero leading coefficients $1,2,\ldots,p-1$, giving one Jordan block.
The bilinearity of the defect $D$ follows from additivity of Adams
operations. Lemma~\ref{C:Jordan_Adams_p} thus gives
\[
 D(E_m,B)\big|_{C_p}=4p b_p(w-1)\ne0.                  \tag{C.3.12}
\]
Nonvanishing follows from the free indecomposable basis, since
$w-1=V_p-V_{p-1}-V_1$. Restriction and forgetting forms commute with
exterior powers. Hence (C.3.12) detects a nonzero Adams
multiplicativity defect in $WG_{\mathrm{sp}}(B_{O_m})$.

For $m=3$ use $SO_3\simeq PGL_2$, valid over every field with $2\ne0$.
Conjugation on the traceless two-by-two matrices preserves the split
quadratic form $-\det$ and has kernel the scalars. It identifies the
smooth connected three-dimensional groups $PGL_2$ and $SO_3$: its
Lie map is injective, and its image is consequently an open subgroup
of the connected target. Let $U=k^2$ be the standard $SL_2$-module and set
\[
 A=H(\operatorname{Sym}^{2p-2}U),\qquad
 B=\operatorname{Sym}^{p-1}U.
\]
The central $\mu_2$ acts trivially on both, so they descend to $PGL_2$.
The invariant alternating form on $U$, tensor-powered $p-1$ times and
restricted by the symmetrizing idempotent, gives $B$ an invariant
nondegenerate symmetric form. The hyperbolic form on $A$ is invariant
by construction. Extend both to $O_3$ along the projection
$g\mapsto\det(g)g$ from $O_3$ onto $SO_3$.

Restrict to the upper-unitriangular constant subgroup $C_p$ of $SL_2$.
On degree-at-most-$2p-2$ polynomials, write $v=t^p-t$. The basis
\[
 1,t,\ldots,t^{p-1};\qquad v,vt,\ldots,vt^{p-2}
\]
is triangular by polynomial degree. Each block is stable under
$t\mapsto t+1$, since $v$ is invariant, and has a single Jordan block
by the same difference calculation. Hence
$\operatorname{Sym}^{2p-2}U|_{C_p}=V_p\oplus V_{p-1}$ and
$B|_{C_p}=V_p$. All Jordan modules are self-dual, so
$A|_{C_p}=2V_p+2V_{p-1}$. As $V_p\otimes V_j=jV_p$,
Lemma~\ref{C:Jordan_Adams_p} gives
\[
 D(A,B)\big|_{C_p}=2p(p-1)(w-1)\ne0.                 \tag{C.3.13}
\]
This proves failure also for $m=3$.

The program \texttt{code/C\_finish\_lowrank\_lambda.py} constructs the
$O_3$ and $O_4$ witnesses for $p=3,5,7$, including their symmetric
forms, actual symmetric-power matrices and tensor-product Jordan
blocks. It independently computes the Adams operations from actual
exterior matrices on every Jordan block. The respective defect
coefficients in front of $w-1$ are $12,40,84$ for $O_3$ and
$36,100,196$ for $O_4$; the largest tensor matrix has dimension $336$.
These computations verify six finite witnesses and every matrix identity
used in their certificates; the preceding argument proves all primes
and all ranks.
\end{proof}

\begin{proposition}[Proved: the rings in ranks zero, one and two over any affine base]
\label{C:BO_lowrank_base}
Let $k$ be any commutative ring with $2\in k^\times$, and put
$W=GW(k)$, the group completion of nondegenerate symmetric projective
$k$-modules under orthogonal sum. Then
\[
 WG_{\mathrm{sp}}(B_{O_0})=W,\qquad
 WG_{\mathrm{sp}}(B_{O_1})=W[\delta]/(\delta^2-1),
\]
where $\delta$ is the sign representation with form $\langle1\rangle$.
For $O_2$, put $d=(\det,\langle-1\rangle)$ and let $u_a$, $a\geq1$,
be the exchanged-weight representation of weights $a,-a$ with Gram
matrix $\left(\begin{smallmatrix}0&1\\1&0\end{smallmatrix}\right)$.
Then the answer is the free $W$-module
\[
 W\,1\ \oplus W\,d\ \oplus\bigoplus_{a\geq1}W\,u_a,
\]
with multiplication
\[
 d^2=1,\quad du_a=u_a,\quad
 u_a u_b=u_{a+b}+u_{|a-b|},\qquad
 u_0=\langle2\rangle(1+d).                             \tag{C.3.14}
\]
For the displayed basis forms,
$\lambda_t(d)=1+dt$ and $\lambda_t(u_a)=1+u_at+dt^2$.
For a coefficient form $(P,b)$, the exterior operations are obtained by
forming the actual finite projective representation $P\otimes u_a$
with its tensor form, taking its exterior power, and decomposing by the
same weight orbits and zero-weight eigenspaces. This coefficient rule
and the orthogonal-sum formula specify every exterior operation,
without any field, noetherian, reducedness or freeness assumption on $k$.
\end{proposition}
\begin{proof}
The standard equivalence between $\mathbb G_m$-comodules and
$\mathbb Z$-graded modules is immediate from the basis $t^a$ of
$k[t,t^{-1}]$: coassociativity makes the coefficient maps orthogonal
idempotents. A finitely generated projective comodule has finitely many
nonzero weight modules, each a projective direct summand. The reflection
in $O_2=\mathbb G_m\rtimes C_2$ exchanges the $a$ and $-a$ modules.
An invariant form pairs only opposite weights, since the Laurent
monomials remain a $k$-basis. On an opposite-weight pair identify the
negative weight with the positive weight by reflection. The pairing
then becomes a symmetric nondegenerate form $b$ on that projective
module; the pair is precisely $(P,b)\otimes u_a$.
On weight zero the idempotents $(1\pm r)/2$ split the reflection;
their images are orthogonal and give the two coefficient forms.
These constructions are inverse equivalences of the corresponding
isometry groupoids under orthogonal sum. Thus group completion gives
the asserted direct sums over $W$. The same argument for $\mu_2$
gives the formula in rank one, and rank zero is the definition.

The matrices in Zibrowius, Lemmas 3.9--3.10
\cite{C-Zibrowius2015}, use only $0,1,-1$; the zero-weight change of
basis additionally has determinant $-2$. They are therefore isometries
over this arbitrary $k$, not just over a field, and yield (C.3.14).
In particular the reflection eigenvectors in the zero-weight
hyperbolic plane have norms $2,-2$, explaining the factor
$\langle2\rangle$. Direct determinants give the stated exterior
polynomials. Taking actual tensors and exterior powers commutes with
isometries and base change, while the weight/eigenspace decomposition
just proved is canonical up to isometry. This proves the coefficient
rule and its completeness.
\end{proof}

\begin{theorem}[Proved: the special-$\lambda$ criterion over every affine base]
\label{C:BO_affine_lambda}
Let $k$ be a commutative ring with $2\in k^\times$, and let $m\geq0$.
Then
\[
 WG_{\mathrm{sp}}(B_{O_m,k})\text{ is special under exterior powers}
 \quad\Longleftrightarrow\quad
       m\leq2\ \text{or}\ k\text{ is a }\mathbb Q\text{-algebra}.       \tag{C.3.15}
\]
There is no noetherian, reducedness, characteristic-purity or
finite-dimensionality hypothesis on $k$.
\end{theorem}
\begin{proof}
We first explain the exact-group input over arbitrary $k$.
Fasel--Haution, Proposition 4.2.2 and the proof of Corollary 4.2.3
\cite{C-FaselHaution2025}, establish the universal special-$\lambda$
identities in exact Grothendieck--Witt groups of $O_r$ and
$O_r\times O_s$ over $\mathbb Z[1/2]$. Their identities transfer to
exact groups of equivariant projective $k$-modules as follows.
For a symmetric $G$-representation $(P,q)$ of rank $r$, its orthogonal
frame torsor $T\to\operatorname{Spec}k$ has a left $G$-action commuting
with its right $O_r$-action. Associate to an $O_r$-representation $V$
over $\mathbb Z[1/2]$ the equivariant projective module
$T\times^{O_r}(V\otimes k)$. This is an exact tensor functor preserving
forms, isometries, Lagrangians, and exterior powers. Exactness follows
on the faithfully flat trivializing cover, where it is scalar extension
of a sequence of finite projective $\mathbb Z[1/2]$-modules; the latter
sequence is split on underlying modules. A pair of frame torsors gives
the identical functor for $O_r\times O_s$.
Thus the universal iterated-exterior and product identities hold in
the exact equivariant group. Their extension from actual objects to
virtual classes uses the pre-$\lambda$ addition formula. Finitely many
open-and-closed rank strata suffice for every pair of projectives, so
constant rank need not be assumed.

For $m\leq2$, every admissible sequence of $O_m$-representations
splits equivariantly. In ranks zero and one this follows from
projectivity and the two sign idempotents; in rank two use the weight
and reflection decomposition in Proposition~\ref{C:BO_lowrank_base}.
On each multiplicity module the quotient is projective, so its sequence
splits over $k$, and these splittings reconstruct an equivariant one.
Any equivariant Lagrangian consequently has an equivariant complement,
and the usual correction by half its Gram map makes that complement
isotropic. Metabolic relations are therefore already actual isometries,
and hyperbolic exact-sequence relations are already orthogonal-sum
isometries. The split and exact groups coincide. The preceding
universal identities prove specialness.

If $k$ is a $\mathbb Q$-algebra, the same conclusion holds for every
$m$. The coordinate coalgebra $\mathbb Q[O_m]$ is a direct sum of
matrix coalgebras, since its finite-dimensional representation category
is split semisimple. After tensoring with $k$, every finite projective
comodule has a finite isotypic decomposition, whose multiplicity modules
are projective direct summands. A sequence therefore splits on each
multiplicity module and hence equivariantly. The preceding comparison
again identifies the split and exact groups and proves specialness.

Conversely assume $m\geq3$ and $k$ is not a $\mathbb Q$-algebra.
Some prime $p$ is not a unit of $k$; since $2$ is a unit, $p>2$.
Choose a maximal ideal containing $p$ and let $F$ be its residue field.
We now lift the negative witnesses to $\mathbb Z[1/2]$, so that their
nonvanishing over $F$ detects nonvanishing over $k$.
For $m\geq4$ take
\[
 E_m,\qquad H(\operatorname{Sym}^{p-1}E_m).
\]
Both are nondegenerate symmetric representations over $\mathbb Z[1/2]$;
no factorial is inverted to form the hyperbolic object. Over $F$ the
symmetric power is self-dual by the form constructed in
Theorem~\ref{C:BO_field_lambda}, so the Green-ring defect is twice
(C.3.12) and is nonzero. For $m=3$ the even-degree representations
$\operatorname{Sym}^{2p-2}(\mathbb Z[1/2]^2)$ and
$\operatorname{Sym}^{p-1}(\mathbb Z[1/2]^2)$ descend to $PGL_2$
over $\mathbb Z[1/2]$. Extend them to $O_3$ as before and take the
hyperbolic forms on both. Over $F$ their defect is twice (C.3.13),
again nonzero. Scalar extension $k\to F$ is a homomorphism of split
pre-$\lambda$-rings, since it preserves isometries, sums, tensors and
exterior powers. A zero defect over $k$ would have zero image over
$F$, contradiction. This proves necessity and the equivalence.
\end{proof}

\begin{proposition}[Proved: explicit rings over every rational affine base]
\label{C:BO_Qalgebra}
Let $k$ be any commutative $\mathbb Q$-algebra. For split $O_m$ choose
the rational simple representations $V_\lambda$ and rational invariant
symmetric forms $q_\lambda$ indexed as in
Proposition~\ref{C:orthogonal_representation_ring}. Then
\[
 WG_{\mathrm{sp}}(B_{O_m,k})
   =\bigoplus_{\lambda\in\operatorname{Irr}(O_{m,\mathbb Q})}
          GW(k)[V_\lambda\otimes k,q_\lambda\otimes k].          \tag{C.3.16}
\]
The tensor coefficients are the rational symmetric multiplicity forms
of Proposition~\ref{C:orthogonal_representation_ring}, extended to $k$.
For exterior operations one forms the actual exterior power of a finite
sum $(V_\lambda\otimes P_\lambda,q_\lambda\otimes b_\lambda)$ and
extracts its projective isotypic multiplicity forms. This gives the
whole ring and all its exterior operations over arbitrary nonfield,
nonreduced and nonnoetherian rational affine bases. It is special by
Theorem~\ref{C:BO_affine_lambda}.
\end{proposition}
\begin{proof}
Use the matrix-coalgebra decomposition from the preceding proof.
Concretely, a finite projective representation has coaction through a
finite direct sum of these coalgebras; dualizing that finite sum gives
a product of matrix algebras over $k$. The matrix idempotents exhibit
each multiplicity module $P_\lambda$ as a finite projective direct
summand and give an actual decomposition
$M=\bigoplus_\lambda V_\lambda\otimes P_\lambda$.
The simple rational $O_m$-modules have scalar endomorphisms and are
symmetric and self-dual, by the opened Zibrowius Proposition 3.17.
Consequently pairings between different components vanish, and the
pairing on one component is uniquely $q_\lambda\otimes b_\lambda$.
Nondegeneracy is equivalent to that of $b_\lambda$. This proves the
isometry classification and its group completion (C.3.16).
Tensor decompositions over $\mathbb Q$ extend to $k$ and tensor the
multiplicity forms with $b_\lambda\otimes b_\mu$. Exterior powers
remain finite projective representations, so the same canonical
isotypic extraction applies. This proves the asserted operations and
that no additional additive relations or objects are omitted.
\end{proof}

\begin{proposition}[Proved reduction: the remaining rank-three module problem contains a wild category]
\label{C:BO3_wild_gate}
Let $k$ be algebraically closed of characteristic $p>2$, and put
$d=p^2+1$ if $p\geq5$, and $d=12$ if $p=3$. There is a fully faithful
exact functor
\[
 F:\operatorname{mod}S_k(2,d)\longrightarrow
             \operatorname{Rep}^{\mathrm{fd}}_k(O_3),
\]
where $S_k(2,d)$ is a Schur algebra of wild representation type.
The functor preserves indecomposables. With
$M^\dagger=M^*\otimes\det^d$, viewed as a degree-$d$ polynomial
$GL_2$-module, for indecomposable $M,N$ one has
\[
 [H(FM)]=[H(FN)]\text{ in }WG_{\mathrm{sp}}(BO_3)
 \quad\Longleftrightarrow\quad N\simeq M\text{ or }N\simeq M^\dagger.
                                                               \tag{C.3.17}
\]
Thus the additive split group retains these wild modules up to the
specified duality. This is a rigorous reduction of the outstanding
classification problem, not an enumeration of the indecomposables
and not a completed answer to the full modular-ring question.
\end{proposition}
\begin{proof}
The coalgebra of homogeneous degree-$d$ polynomials in the four matrix
coordinates has dual $S_k(2,d)$. Finite-dimensional modules over its
dual are exactly its finite-dimensional comodules, hence homogeneous
polynomial $GL_2$-representations of degree $d$. This equivalence,
including the indicated scalar weight, is stated in
Doty--Erdmann--Martin--Nakano, Section 1.3, printed p.140
\cite{C-DotyErdmannMartinNakano1999}.
Their Proposition in Section 3.4, printed pp.151--152, says that
$S_k(2,d)$ is wild for $p\geq5,d\geq p^2$, and for $p=3,d\geq12$.
The precise cases are confirmed by
Hemmer--Kujawa--Nakano, Theorem 1.2.1(b)--(d)
\cite{C-HemmerKujawaNakano2006}. The original authors' correction for
$S(2,11)$ in characteristic two does not affect either chosen case.

On every homogeneous degree-$d$ module scalar matrices act by $t^d$.
Therefore a linear map between two such modules commuting with $SL_2$
also commutes with scalar matrices. Over algebraically closed $k$ every
$g\in GL_2(k)$ is a scalar times an element of $SL_2(k)$, by choosing
a square root of $\det(g)$. Such a map is consequently $GL_2$-equivariant;
the equality of rational morphisms follows from density of $k$-points
in the smooth reduced group. Restriction to $SL_2$ is fully faithful.
Since $d$ is even, the central $\mu_2$ acts trivially, so these modules
descend to $PGL_2\simeq SO_3$. Extend to $O_3$ along
$g\mapsto\det(g)g$. All these functors are exact on vector spaces and
fully faithful, proving the assertion about $F$. Its equality of
endomorphism rings proves preservation of indecomposables.

For clarity $M^\dagger$ is again polynomial of degree $d$:
the inverse matrix formula has denominator $\det(g)$, so the degree-$d$
matrix coefficients of the dual become polynomial after multiplication
by $\det(g)^d$, with total degree $d$. On $SL_2$ that determinant
twist is trivial, giving $F(M^\dagger)\simeq F(M)^*$.
If the two classes in (C.3.17) agree, forgetting forms gives
$[FM]+[(FM)^*]=[FN]+[(FN)^*]$ in the split Green group.
Krull--Schmidt and the free indecomposable basis give
$FN\simeq FM$ or $FN\simeq(FM)^*$. Full faithfulness gives the two
asserted possibilities. Conversely an isomorphism of modules induces
an isometry of their hyperbolic forms, and $H(V)\simeq H(V^*)$ by
exchanging its two summands. This proves the reverse implication and
specifies exactly the ambiguity retained by hyperbolic passage.
\end{proof}
